\section{Justificación}

En la actualidad, un número creciente de personas manifiesta preocupación por el
medio ambiente y aspira a consumir de manera más responsable. Sin embargo, el
mercado no facilita esta tarea. En primera instancia, se encuentra el
greenwashing, práctica mediante la cual las empresas emplean un lenguaje
vagamente ecológico para aparentar sostenibilidad sin serlo realmente. Mensajes
sobre ``compromiso con el planeta'' o ``materiales reciclados'' frecuentemente
encubren productos que siguen siendo difíciles de reparar y de vida útil
reducida. Esto genera desconfianza: el consumidor no sabe en quién confiar
\parencite[p.~9]{holloway2025}.

A esto se suma la sobrecarga de información. Entre datos técnicos complejos,
especificaciones confusas y campañas de marketing, resulta muy difícil para el
consumidor determinar qué producto es realmente duradero. La ausencia de
estándares claros sobre reparabilidad o durabilidad hace que, aunque la
intención de consumir de forma sostenible exista, la falta de herramientas para
hacerlo termine perpetuando el mismo ciclo: adquirir, usar brevemente,
reemplazar \parencite[p.~9]{holloway2025}.

Frente a esta problemática, la economía circular surge como un marco regulatorio
orientado a contrarrestar precisamente lo que este ciclo de consumo acelerado
genera: el desperdicio desmedido de recursos. En México, aunque esta idea de
economía circular es relativamente nueva, en años recientes ha incrementado el
interés por adoptarla mediante leyes que aún están en discusión. Sin embargo,
existen normas oficiales que, aunque no es su objetivo principal, encaminan al
país a su adopción \parencite[p.~220]{munoz2023}:

\begin{itemize}
    \item \textbf{NOM-052-SEMARNAT-2005:} Establece las características, el
    procedimiento de identificación, clasificación y los listados de los
    residuos peligrosos.
    \item \textbf{NOM-161-SEMARNAT-2011:} Establece los criterios para
    clasificar los Residuos de Manejo Especial y determinar cuáles están sujetos
    a Plan de Manejo.
\end{itemize}

Si bien estas normas representan un avance en la gestión de residuos, buscando
reducir su impacto ambiental, su enfoque se limita al manejo posterior de los
mismos, sin abordar una de las principales causas que los genera: la
obsolescencia programada. 

Es precisamente en este vacío donde se situó el enfoque central de la
investigación. En la actualidad, las empresas comercializan sus productos con
acceso asimétrico a la información: mientras los fabricantes conocen en detalle
la calidad, los riesgos y las limitaciones de sus productos, esta información
raramente es comunicada al consumidor de forma completa o comprensible. Esto
constituye una práctica que perjudica directamente al consumidor, quien toma
decisiones de compra sin contar con los elementos necesarios para evaluarlas con
criterio. Ante esto, la aplicación buscó reducir esta brecha
informativa centralizando, de manera accesible, datos relevantes sobre
durabilidad, reparabilidad, prácticas corporativas de los dispositivos
electrónicos. De esta forma, no solo se combate la desinformación que favorece a
la obsolescencia programada, sino que también se le otorga al consumidor una
herramienta para identificar el greenwashing y navegar en la información con mayor
claridad, fomentando así decisiones de consumo más conscientes y responsables. 

A fin de garantizar la viabilidad del proyecto, se establecieron una serie de
delimitaciones que permitieron mantener un enfoque claro acorde con las
capacidades actuales dentro del contexto de la materia. La investigación se
restringió a productos electrónicos, específicamente smartphones y laptops. En
ambas categorías se abordaron las marcas Apple, Samsung, Huawei y Xiaomi, dado
que representan mayor presencia en el mercado; además, constituyen un sector
donde la información sobre reparaciones suele ser especialmente restrictiva. 

La recolección de datos se limitó a información disponible dentro de
México, llevando a cabo un estudio transversal que permitió analizar el estado
actual de la problemática en el país. Debido a que un análisis de alcance global
implicaría una cantidad de tiempo que exceden los límites del presente proyecto
se eligió la anterior delimitación geográfica. Asimismo, la cantidad de
información manejada se acotó, priorizando aquellos datos más relevantes
para el consumidor, con el fin de garantizar que el diseño se mantenga dentro de
las capacidades técnicas disponibles. 
